\section{What does it mean to forecast inductively?}
\label{sec:inductive-forecasting}

Forecasting is a useful testbed for inductive reasoning because it makes
otherwise hidden choices observable. A forecaster must decide which evidence
matters, which population provides the relevant comparison, how strongly that
evidence should affect an outcome, and when direct observations should override
an existing belief. Each of these choices implies a different numerical
forecast. An accurate prediction alone does not tell us whether the forecaster
made those choices for the right reasons.

This distinction matters because a forecast is rarely the last question. We often
want to know what evidence should change the prediction, whether the same
evidence should matter differently in another population, and what new
observation should trigger revision. A model that reaches the right number
through a familiar base rate, a topical cue, or some other surface regularity may
still fail under exactly these changes. Average accuracy can therefore conceal an
important weakness: two models can give the same forecast while relying on very
different relationships between evidence and outcomes.

We use \emph{inductive forecasting} to describe the stronger behavior we care
about here. An inductive forecaster should not only identify whether evidence is
relevant, but also infer the relationship that determines how that evidence
should affect the outcome. When context changes which comparison is appropriate,
the forecast should change with it. When direct evidence becomes available, that
evidence should be able to revise the earlier contextual inference. And when the
same underlying relationship appears in an unfamiliar setting, the forecaster
should be able to use it without depending on the original surface form.

\begin{figure}[!t]
\centering
\includegraphics[width=\linewidth]{exp1_protocol_four_panels.pdf}
\caption{\textbf{When should a forecast change?} An agent forecasts an
unresolved question, then receives one fictional update with search off. Some
updates support the outcome, some oppose it, and some concern the same topic but
do not bear on the outcome. The labels remain hidden; we ask whether the model
moves when the news matters and whether it moves the right way.}
\label{fig:exp1-pipeline}
\end{figure}

These requirements extend beyond forecasting. A persuasion model should respond
not only to a message, but to who delivers it. A coordination model should
distinguish shared incentives from what participants know about one another. A
policy model should distinguish similar language from the institutional pathway
through which an intervention takes effect. In each case, the conditional
relationship is often the object of interest: who trusts whom, which population
responds to which signal, or which mechanism connects an action to an outcome.
Matching an average outcome while missing these dependencies can produce a
convincing static prediction and a misleading response when the setting changes
\citep{park2023generative,aher2023simulate,argyle2023out,bisbee2024synthetic}.

This suggests a stronger way to evaluate forecasting systems. Rather than asking
only whether a model reaches the correct number, we can ask whether its forecast
changes when the relevant evidence, comparison, or relationship changes. We
therefore proceed through increasingly demanding tests: first whether models
distinguish relevant evidence from topical information, then whether they infer
how strongly evidence should matter from context, and finally whether those
relationships can be revised and carried into new settings.
